% 原书 p.15–p.18（PDF p.23–p.26）
% >>> p023
\section{拉格朗日子流形（Lagrangian Submanifold）}

\begin{keytakeaways}
\begin{itemize}
  \item 几何图像：拉格朗日子流形是辛形式在其上恒为零、且维数恰为流形一半的子流形；
        它是辛几何中「最看不见、却最丰富」的对象。
  \item 数学工具：$\omega|_Y=0$ 与 $\dim Y=\frac12\dim M$、图像构造
        $\mathrm{graph}(\varphi)\subset(M\times M,\pi_1^*\omega-\pi_2^*\omega)$、
        以及余切丛中的闭 1-形式图像 $\Gamma_\alpha$。
  \item 本节结论：法丛同构 $NL\cong T^*L$（命题 1.6.1），由此得到自相交数与欧拉示性数的
        关系 (1.14)；特别地 $(\R^4,\omega_{\mathrm{std}})$ 中唯一的闭可定向拉格朗日曲面是环面。
  \item 刚性对比：拉格朗日子流形处处存在（余切丛中 $df$ 的图像），但紧的\emph{恰当}拉格朗日
        子流形在 $(\R^{2n},\omega_{\mathrm{std}})$ 中不存在（Gromov，定理 1.6.1）。
\end{itemize}
\end{keytakeaways}

回忆 $M$ 上的辛结构 $\omega$ 是非退化的（参见定义 1.1.1 中的条件 (i)）。这\emph{并不}意味着
它限制到子流形 $N\subset M$ 上仍然非退化。这里有一个平凡的例子：
$Y_1=\mathrm{span}\langle x_1\rangle\subset\R^{2n}$ 且 $\omega_{\mathrm{std}}|_N=0$。类似地，
对任意 $1\le k\le n$，可以考虑 $Y_k=\mathrm{span}\langle x_1,\cdots,x_k\rangle$。不仅
$\omega_{\mathrm{std}}$ 在其上不再非退化，而且实际上 $\omega_{\mathrm{std}}|_N=0$。沿这个方向
继续玩下去，这种「恒为零」第一次失效是在 $y_i$ 加入的时候。换言之，子空间
\[
  Y_n=\mathrm{span}\langle x_1,\cdots,x_n\rangle
\]
是 $\omega$ 恒为零的最大可能子空间，注意 $\dim Y_n=\frac12\dim\R^{2n}$。这个子空间 $Y_n$
代表了辛流形 $(M,\omega)$ 中一大族子流形，它们在辛几何的各个课题中扮演着重要角色。
% [Tutor Note: 原文疑误——本段先用 $Y_1$ 定义子空间，随后却写成「$\omega_{\mathrm{std}}|_N=0$」
% 与「$Y_k$ 上……$\omega_{\mathrm{std}}|_N=0$」，其中 $N$ 并未定义；按上下文应分别为
% $\omega_{\mathrm{std}}|_{Y_1}=0$ 与 $\omega_{\mathrm{std}}|_{Y_k}=0$。译文保留原文记号并在此说明，
% 详见 review_ch01.md。几何上这很直观：$\mathrm{span}\langle x_1,\cdots,x_k\rangle$ 张成的坐标
% 方向两两都是「同类的」，故 $\dd x_i\wedge\dd y_i$ 在其上为零；只有当某个 $y_i$ 与对应的 $x_i$
% 一起出现时，$\omega$ 才可能不为零。]

\begin{wdefn}{1.6.1}
子流形 $Y\subset(M,\omega)$ 称为\textbf{拉格朗日的}（Lagrangian），若
(i) $\omega|_Y=0$，且 (ii) $\dim(Y)=\frac12\dim(M)$。
\end{wdefn}

\begin{wrem}{1.6.1}
注意定义 1.6.1 中的条件 (ii) 正是使定义 1.6.1 中条件 (i) 成立的最大维数（参见本节开头的
动机性例子以及定理 1.2.1）。
\end{wrem}

乍看之下，拉格朗日子流形似乎毫无用处，因为辛结构根本不「看见」它。然而，拉格朗日子流形其实
非常丰富，下面列出其中一些重要的例子。

\begin{wexam}{1.6.1}
辛曲面 $(\Sigma,\omega_{\mathrm{area}})$ 中的任何曲线 $Y$ 都是拉格朗日的，因为对 $Y$ 上的任意
切向量 $v$ 都有 $\omega_{\mathrm{area}}(v,v)=0$。\hfill$//$
\end{wexam}

\begin{wexam}{1.6.2}
对任意 $\varphi\in\Symp(M,\omega)$，$\mathrm{graph}(\varphi):=\{(p,\varphi(p))\mid p\in M\}$
是 $(M\times M,\pi_1^*\omega-\pi_2^*\omega)$ 中的拉格朗日子流形。事实上，在任意点
$(p,\varphi(p))\in\mathrm{graph}(\varphi)$ 处，
\[
  T_{(p,\varphi(p))}\mathrm{graph}(\varphi)
  =\{(v,\varphi_*(v))\mid v\in T_pM\}.
\]
于是
$(\pi_1^*\omega-\pi_2^*\omega)\big((v,\varphi_*(v)),(w,\varphi_*(w))\big)=0$，因为
$\varphi^*\omega=\omega$。我们也可以说其逆同样成立：若 $\mathrm{graph}(\varphi)$ 是
$(M\times M,\pi_1^*\omega-\pi_2^*\omega)$ 中的拉格朗日子流形，则 $\varphi$ 是
% >>> p024
辛同胚。这提供了一个检验微分同胚 $\varphi$ 是否为辛同胚的几何判别准则
（参见问题 1.5.2 的 (i)）。\hfill$//$
\end{wexam}
% [Tutor Note: 逆命题的证明要点：$\mathrm{graph}(\varphi)$ 拉格朗日等价于
% $0=(\pi_1^*\omega-\pi_2^*\omega)|_{\mathrm{graph}(\varphi)}$，而后者在向量
% $(v,\varphi_*v)$ 上取值恰为 $\omega(v,v)-\omega(\varphi_*v,\varphi_*v)$ 类型的表达式；
% 把它写成 $\varphi^*\omega-\omega$ 在 $\mathrm{graph}(\varphi)\cong M$ 上的拉回，即得
% $\varphi^*\omega=\omega$。这里「图像投影到第一个因子是同构」这一事实是关键，
% 对一般拉格朗日子流形（未必是图像）不能这样论证。]

\begin{wrem}{1.6.2}
作为例 1.6.2 的特例，由于 $\mathbf{1}\in\Symp(M,\omega)$，对角线
$\Delta_M:=\{(x,x)\mid x\in M\}$ 是 $(M\times M,\pi_1^*\omega-\pi_2^*\omega)$ 中的拉格朗日
子流形。注意这一点对辛结构的选择是敏感的。例如，取 $M=\R^{2n}$ 并考虑 (1.4) 中定义的积辛结构，
则反对角线 $\overline{\Delta}_M=\{(x,-x)\mid x\in\R^{2n}\}$ 是拉格朗日的，而 $\Delta_M$
不是拉格朗日的。
\end{wrem}
% [Tutor Note: 原文疑误——须把两种辛形式分清。在\emph{差}形式
% $\pi_1^*\omega-\pi_2^*\omega$ 下，$\mathrm{graph}(\varphi)$ 拉格朗日当且仅当
% $\varphi^*\omega=\omega$（即 $\varphi\in\Symp$），故对角线 $\Delta_M$（$\varphi=\mathbf{1}$）
% 与反对角线 $\overline{\Delta}_M$（$\varphi=-\mathbf{1}$，而 $-\mathbf{1}$ 也是辛的）都是
% 拉格朗日的。在\emph{和}形式 $\pi_1^*\omega+\pi_2^*\omega$ 下，
% $\mathrm{graph}(\varphi)$ 拉格朗日当且仅当 $\varphi^*\omega=-\omega$（$\varphi$ 反辛），
% 此时 $\Delta_M$ 确实不拉格朗日，但拉格朗日的是复共轭 $z\mapsto\bar z$ 这类「反全纯对角线」，
% 而非 $\{(x,-x)\}$：后者的切向量 $(v,-v)$ 满足
% $(\pi_1^*\omega+\pi_2^*\omega)((v,-v),(w,-w))=2\omega(v,w)\neq0$。因此原文此句无论按 (1.4)
% 的和形式还是按差形式解读都不成立，疑为作者把「反对角线」与「复共轭图像」或把和形式与差形式
% 混用。译文保留原文写法，详见 review_ch01.md。]

\begin{wexam}{1.6.3}
在余切丛 $(T^*Q,\omega_{\mathrm{can}})$ 中，任意 1-形式 $\alpha\in\Omega^1(Q)$ 的图像是
$T^*Q$ 的一个子流形，记作 $\Gamma_\alpha$。则 $\Gamma_\alpha$ 是
$(T^*Q,\omega_{\mathrm{can}})$ 中的拉格朗日子流形，当且仅当 $\alpha$ 是闭的。为证明这一点，
回忆 $\omega_{\mathrm{can}}=\dd\lambda_{\mathrm{can}}$，其中 $\lambda_{\mathrm{can}}$ 是由 (1.3)
定义的典范原形式。注意 $\lambda_{\mathrm{can}}$ 满足 $\alpha^*\lambda_{\mathrm{can}}=\alpha$。
于是对内含映射 $i:\Gamma_\alpha\to T^*Q$，我们有
\[
  \omega_{\mathrm{can}}|_{\Gamma_\alpha}=i^*\omega_{\mathrm{can}}
  =i^*(\dd\lambda_{\mathrm{can}})=\dd(i^*\lambda_{\mathrm{can}})
  =\dd(\alpha^*\lambda_{\mathrm{can}})=\dd\alpha .
\]
因此我们得到所需的结论。由于这一结果，人们可以通过考虑 $\alpha=\dd f$（$f:Q\to\R$ 光滑）
很容易地在 $(T^*Q,\omega_{\mathrm{can}})$ 中构造出大量拉格朗日子流形（参见定义 1.6.2）。\hfill$//$
\end{wexam}

\begin{wexer}{1.6.1}
证明上面定义的 $\Gamma_\alpha$ 是 $(T^*Q,\omega_{\mathrm{can}}=\dd\lambda_{\mathrm{can}})$ 中的
\emph{恰当}拉格朗日子流形（见定义 1.6.2），当且仅当 $\alpha$ 是恰当 1-形式。
\end{wexer}

\begin{wrem}{1.6.3}
一个拉格朗日子流形 $Y$ 可以带许多奇点地浸入 $M$，它也可以是非紧的或不可定向的。
\end{wrem}

沿相反的方向（与拉格朗日子流形的丰富性相对），拉格朗日子流形本身也受到一些拓扑约束。
下面就是一条。

\begin{wprop}{1.6.1}
对任意拉格朗日子流形 $L$，其法丛 $NL$ 同构于其余切丛 $T^*L$。
\end{wprop}

\begin{wproof}
回忆对任意 $p\in L$，由丛的短正合列 $0\to TL\to T_LM\to NL\to0$（其中 $T_LM$ 表示限制在
$L$ 上的丛 $TM$）给出的 $NL$ 的定义，我们有
\[
  N_pL\cong(T_LM)_p/T_pL=T_pM/T_pL .
\]
同时，对任意 $[v]\in T_pM/T_pL$，考虑如下映射
\[
  \sigma_p:T_pM/T_pL\longrightarrow T_p^*L
  \qquad\text{由}\qquad [v]\mapsto\omega_p(v,\cdot)\ \text{定义}.
\]
注意 $\sigma_p$ 是良定义的，因为当 $v\in T_pL$ 时 $\sigma_p([v])=0$（由于 $L$ 是拉格朗日的）。
由维数理由，只需证明 $\sigma_p$ 是单射。事实上，若 $\omega_p(v,\cdot)$ 在 $T_p^*L$ 上为 $0$，
为了不违背 $\omega$ 的非退化条件，向量 $v$ 必须落在 $T_pL$ 之中。否则，子空间
$\mathrm{span}\langle v,T_pL\rangle$ 就给出 $T_pM$ 的一个 $(n+1)$ 维子空间，
% >>> p025
在其中 $\omega_p$ 恒为零，这与 $L$ 是拉格朗日的定义相矛盾（见注 1.6.1）。因此
$[v]=0\in T_pM/T_pL$。于是我们断定 $\sigma_p$ 是同构。
\end{wproof}
% [Tutor Note: 原文疑误——「若 $\omega_p(v,\cdot)$ 在 $T_p^*L$ 上为 0」应为「在 $T_pL$ 上为 0」：
% $\omega_p(v,\cdot)$ 作为余向量须作用于切向量，作用对象是 $T_pL$ 而不是余切空间 $T_p^*L$。
% 原文此处多写了一个星号，译文保留原文写法，详见 review_ch01.md。
% 另外，$\sigma_p$ 的「单射」论证实质是：若 $v\notin T_pL$，则
% $\dim\,\mathrm{span}\langle v,T_pL\rangle=n+1$ 且 $\omega$ 在其上恒为零，这与
% 「辛向量空间的极大迷向子空间维数为 $n$」矛盾；而「极大迷向子空间维数为 $n$」由命题 1.1.2
% （$k=1$ 的情形）或线性代数中斜对称形式的秩定理保证。]

命题 1.6.1 直接推出，当拉格朗日子流形 $L$ 闭且可定向时，有如下拓扑障碍。记 $e(E)$ 为向量丛
$E$ 的欧拉类，则
\begin{equation}
  [L]\cdot[L]=\int_L e(NL)=-\int_L e(TL)=-\chi(L),
  \tag{1.14}
\end{equation}
其中 $[L]\cdot[L]$ 是 $L$ 的自相交数，而 $\chi(L)$ 是 $L$ 的欧拉示性数。在
$(\R^4,\omega_{\mathrm{std}})$ 中，任何闭的（拉格朗日）曲面 $L=\Sigma_g$ 都可以被扰动而离开
自身，故 $[L]\cdot[L]=0$。于是 (1.14) 蕴涵其欧拉示性数必为 $0$。由于任何闭可定向曲面都由其
亏格分类，且 $\chi(L)=\chi(\Sigma_g)=2-2g$，我们知道 $\chi(L)=0$，从而 $g=1$。因此，
$(\R^4,\omega_{\mathrm{std}})$ 中唯一的闭可定向拉格朗日曲面就是二维环面！有趣的是，对不可定向的
拉格朗日子流形，例如 Klein 瓶，情形会更加复杂 [167,206]。

到目前为止的讨论表明，研究单个拉格朗日子流形有时并不那么有趣。在许多情形下，研究某一族
拉格朗日子流形才更引人入胜。

\begin{wexam}{1.6.4}
考虑 $(\R^{2n},\omega_{\mathrm{std}})$ 与 $L_*=\mathrm{span}\langle x_1,\cdots,x_n\rangle$。
用 $A\cdot L_*$（$A\in\mathrm{Sp}(2n)$）可以得到许多其他的拉格朗日线性子空间。事实上，
\[
  \Lambda_n:=\{\text{所有 }(\R^{2n},\omega_{\mathrm{std}})\text{ 的拉格朗日线性子空间}\}
\]
被 $\mathrm{Sp}(2n)$ 传递地作用。下面的练习展示了一个众所周知的结果。\hfill$//$
\end{wexam}

\begin{wexer}{1.6.2}
假设 $\Lambda_n$ 具有流形结构（它被称为\textbf{拉格朗日 Grassmann 流形}），证明
\[
  \dim_{\R}\Lambda_n=\frac{n(n+1)}{2}.
\]
\emph{提示}：应用代数结果 $G\cdot p\simeq G/G_p$，其中 $G\cdot p$ 是轨道集而 $G_p$ 是 $p$ 的
稳定化子，取点 $p=L_*$ 与群 $G=\mathrm{Sp}(2n)$（于是 $G\cdot p=\Lambda_n$）。
\end{wexer}

为了结束本节，让我们指出一族特殊的拉格朗日子流形，它们对检验许多刚性结果很有用
（例如见第 10 章 10.1 节）。

\begin{wdefn}{1.6.2}
在恰当辛流形 $(M,\dd\lambda)$ 中，拉格朗日子流形 $L\subset(M,\dd\lambda)$ 称为
\textbf{恰当的}（exact），若 $\lambda|_L=\dd f$，其中 $f:L\to\R$ 为某个光滑函数。
\end{wdefn}

由例 1.6.3 与练习 1.6.1，我们已经看到 $(T^*Q,\omega_{\mathrm{can}})$ 中含有大量由恰当
1-形式 $\alpha=\dd f$ 的图像给出的恰当拉格朗日子流形。然而，下面这个同样归功于 Gromov 的定理，
与欧氏空间的情形形成了尖锐的对比。

\begin{wthm}{1.6.1}（Gromov [99]）
在 $(\R^{2n},\omega_{\mathrm{std}})$ 中不存在紧的恰当拉格朗日子流形。
\end{wthm}

% >>> p026
本书不会证明这一结果。作为替代，我们将在第 9 章 9.7 节中，通过应用一个源自 $J$-全纯曲线的
高级工具，来证明 $(\CP^n,\omega_{\mathrm{FS}})$ 中的一个类似结果。
